%This is a template file for use of iopjournal.cls

\documentclass{iopjournal}
\usepackage{array}
\usepackage{longtable}
\usepackage{float}
\usepackage{amsmath}
\usepackage[utf8]{inputenc}
\usepackage{siunitx}

\begin{document}

\articletype{Research Letter} %	 e.g. Paper, Letter, Topical Review..

\title{Waste Data: Operational Dairy Records for Pollution Accounting}

% Alt titles
% Waste Data: Operational dairy records reveal missing nitrogen and fixable reporting gaps
% Dirty Data: Operational dairy records reveal missing nitrogen and fixable reporting gaps
% Waste Data: Assessing Operational Dairy Records for Pollution Accounting
% Waste Data, Un-wasted: Operational Dairy Records for Pollution Accounting
% Waste data, not wasted data: Operational Dairy Records for Pollution Accounting
% Operational Dairy Records Reveal Unaccounted-for Manure and Transportation Emissions

\author{I. Avery Bick$^{1,2}$*\orcid{0000-0002-8007-9649}, Daly Wettermark$^3$\orcid{0000-0001-7151-2765}, Gina Hervey$^{2}$\orcid{0009-0002-3409-5135}, Elena Eneva$^{1}$\orcid{0009-0001-8344-9609}, Daniel E. Ho$^{1}$\orcid{0000-0002-2195-5469}}

\affil{$^1$Regulation, Evaluation and Governance Lab, Stanford University, Stanford, CA, USA}

\affil{$^2$Climate and Energy Policy Program, Stanford University, Stanford, CA, USA}

\affil{$^3$Department of Civil and Environmental Engineering, Stanford University, Stanford, CA, USA}

\affil{$^*$Author to whom any correspondence should be addressed.}

\email{averybick@stanford.edu}

\keywords{dairy, manure, agricultural policy, water quality, greenhouse gas inventory}

\begin{abstract}
Agriculture has proven to be one of the most vexing sectors for pollution accounting. The diversity of agricultural facilities, non-point source nature of emissions, and lack of available data have stymied facility-level accounting by governments, nonprofits, and researchers alike. Here, we explore one novel source of potential data to benchmark environmental impacts of agriculture: operational data from concentrated animal feeding facilities, which have come under increased scrutiny due to nitrogen discharges affecting groundwater resources and significant greenhouse gas emissions. First, we process annual reports and manure transportation manifests for 795 dairy operations within California's Central Valley Regional Water Quality Control Board in 2024 and extract relevant information using an agentic approach and human-in-the-loop validation. Second, we perform a nitrogen balance across dairies in the region, estimating that 41.7 million lbs of nitrogen was unaccounted-for in 2024. We find that dairies with less capacity to remove nitrogen via crops tend to have higher fractions of unaccounted-for nitrogen, raising the possibility of undocumented over-application, which could threaten already stressed groundwater resources. Third, we model manure trucking routes based on transport manifests, estimating that heavy trucks traveled at least 2.43 million miles and consumed at least 378,000 gallons of diesel in the process of manure transportation in 2024, disproportionately driving through low-income and unincorporated communities. Lastly, we show that California's greenhouse gas inventory rests on brittle assumptions regarding manure management and explore the potential of operational data to update them, and demonstrate how this is limited by both noncompliance with record-keeping requirements and the reporting design itself. Overall, we demonstrate that previously difficult-to-digest operational datasets can be increasingly leveraged to identify potential strategic misreporting of discharges and improve pollution accounting for industrial animal agriculture.
\end{abstract}

\section{Introduction}

The Central Valley contains among the most nitrate-contaminated groundwater aquifers \cite{Balazs_2011, USGS_2010} in the United States, as well as the highest annual average concentration of atmospheric ammonia \cite{Wang_2021} and particulate matter ($\text{PM}_{2.5}$) \cite{Chen_2020}. Concentrated animal feeding operations (CAFOs), such as dairies, have long been acknowledged to substantially contribute to nitrate contamination of groundwater \cite{Harter_Lund_2012}, odor issues \cite{Rabaud_2003, Akter_2020, Lu_2008}, and ammonia \cite{Miller_2015} and $\text{PM}_{2.5}$ air pollution, primarily as a result of manure management. The Central Valley includes some of the lowest-income communities in California who experience increased respiratory illness rates \cite{Zarate_2024}, and have relatively limited resources for treating or replacing contaminated water resources \cite{Harter_Lund_2012}. Dairy facilities also have large impacts on California's greenhouse gas emissions inventories; dairy manure and dairy cattle enteric fermentation combined accounted for an estimated 44.7\% of California's methane emissions in 2020 \cite{El_Mashad_2023, CARB_GHG_2022}. 

%Despite being in ``serious non-attainment'' status for $\text{PM}_{2.5}$ \cite{Valley_Air_2024} and ``extreme non-attainment'' for ozone \cite{Valley_Air_Ozone_2024} under the Clean Air Act, the San Joaquin Valley Air District has a myriad of issues with the emissions estimates underlying its dairy permitting process. The District only includes $\text{PM}_{2.5}$ emissions from cattle housing, excluding emissions from manure handling and land application \cite{SJVAPCD_Permit_2025}, and emissions factors may be based on unpublished research, such as for volatile organic compounds \cite{ValleyAir_2012}, an ozone precursor. These emissions modeling issues are compounded by public data gaps. In the District, evidence of facility-level compliance and exemptions are kept on-site at dairies \cite{Valley_Air_Rule_2020}. There is also great uncertainty regarding manure management and manure endpoints \cite{Robertson_2026}, which may sizably impact these facilities' greenhouse gas emissions estimates.

Reliable and open reporting data is a critical tool for establishing environmental compliance, allowing public oversight, and informing future policy \cite{Giles_2022}. Governments collect vast quantities of data from facilities to monitor emissions and discharges and enforce environmental compliance \cite{Gray_Shimshack_2011, Giles_2022}. The 2018 Foundations for Evidence-Based Policymaking Act increased attention on the potential of such data for analyses and policy development \cite{US_Congress_2019}. However, in practice such data can be limited by reporting design \cite{Bennear_2008} as well as by intermittent oversight and compliance \cite{Robertson_2026}. In the agricultural sector, facility-scale operational data can often be unreliable, incomplete, or simply unavailable \cite{Marklein_2021, Miller_2020, ELF_v_SWRCB_2023, Handan-Nader_Ho_2019, Robertson_2026}. Accounting efforts such as the California Air Resources Board (CARB) Greenhouse Gas (GHG) Inventory may thus rely on generalized approaches and broad assumptions on where manure ultimately goes \cite{CARB_2000_2023_GHG_Methods}. 

For instance, manure management system coverage data in the 2025 edition of CARB's GHG Inventory program \cite{CARB_2000-2023_GHG} are based on federal reports and farm surveys from 2007 \cite{EPA_Inventory_Annexes_2013}. Outdated survey data does not reflect the increase in dairy sizes, which nearly doubled between 2002 and 2017 \cite{USDA_2020}. Larger dairies are more likely to use anaerobic lagoons \cite{EPA_Inventory_Annexes_2013} and digesters \cite{Cowley_2018} which can generate significant methane emissions from direct off-gassing \cite{Marklein_2021} and leakage \cite{Valdez_2026}, respectively. Smaller dairies tend to rely on more aerobic solid storage and composting systems \cite{EPA_Inventory_Annexes_2013, Owen_2015}. Manure management systems have drastic implications for methane emissions, which can span three orders of magnitude depending on the manure management practices used \cite{Marklein_2021, El_Mashad_2023, Naranjo_2020}. Overall, more temporally relevant information could improve the representativeness of CARB's current methodologies for estimating emissions from the dairy sector.

California's Central Valley Regional Water Quality Control Board (the ``Water Board'') regulates waste discharges that affect water quality in the Central Valley, a significant dairy producing region that contains almost 90 percent of the State's cattle population \cite{SWRCB_2026_Order}. The region faces significant groundwater nitrate pollution due to agricultural applications of synthetic fertilizers and dairy manure \cite{Harter_Lund_2012, SWRCB_2026_Order}, which can lead to a myriad of adverse health outcomes, including cancers and methemoglobinemia (``blue-baby syndrome'') \cite{Ward_2018}, and increased treatment costs for municipal drinking water systems \cite{Jensen_2011}. In an effort ``to confirm that the discharges are effectively controlled by management practices'', the Water Board's Monitoring and Reporting Program mandates that dairies must submit annual reports \cite{CVRWQCB_2013_Order} that detail herd size, the amount of manure and wastewater generated and applied to dairy-owned fields \cite{Robertson_2026}. Dairies are required to limit the ratio of applied nitrogen to crop-removed nitrogen (the ``applied-to-removed" ratio) on their fields to 1.4 \cite{CVRWQCB_2013_Order}. These limits create pressure on dairies, especially those with limited cropland, to reduce application. However, off-site transportation of manure can be a major cost for dairies that are trying to meet application limits \cite{USDA_2003}. Thus, annual reports for Central Valley dairies are required to include manure transportation manifests to document exports and ensure compliance \cite{CVRWQCB_2013_Order}. 

While the Water Board dairy annual reports have been criticized due to data quality issues and incomplete accounting of manure and nutrients \cite{Robertson_2026, SWRCB_2026_Order}, such facility-scale operational data has the potential to refine GHG and nitrogen pollution accounting for California's dairies. Towards these goals, we collect, clean, and structure data from 795 dairy facility annual reports containing 1,264 manure and wastewater transport records from within the Water Board's jurisdiction, which includes the majority of the State's dairies. From this, we create a dataset of self-reported herd characteristics, manure and wastewater generation, land application, and exports for Central Valley dairies. This allows us to perform facility-level and Water Board-level nitrogen accounting, from which we identify patterns of unaccounted-for nitrogen reporting among dairies with less capacity for crop uptake of nitrogen. We then utilize manure transportation manifests and OpenStreetMaps road network data \cite{OSM_California_2025} to model trucking routes between dairies and destinations and estimate transportation emissions and exposures at the community-level \cite{Census_Places}. Lastly, we explore applicability of operational data outside of the Water Board setting, comparing the manure management methods assumed by CARB's Greenhouse Gas (GHG) Emissions Inventory \cite{CARB_2000_2023_GHG_Methods} against those we observe in annual reports. Overall, we demonstrate that operational data can be used for pollution accounting and detection of potential misreporting in agricultural settings and that simple interventions could improve reporting quality design.

\begin{figure}[H]
 \centering
        \includegraphics[width=1.0\textwidth]{Figures/Figure_1_Data_Flow.jpg}
 \caption{Dairy operational data is submitted in many formats, necessitating a data extraction and validation pipeline. a) Data flow for processing 2024 Central Valley dairy operational data: 4 formats of annual reports and 14 formats of manure and wastewater transportation manifests. b) An example dairy annual report in a non-standard format. c) A hand-written manure transport manifest from the same dairy.}
\label{DataFlow}
\end{figure}

\section{Methods}
\subsection{Characterizing Dairy Operational Data}

The dataset consists of 859 dairy annual reports from the Water Board for the year 2024 \cite{Robertson_2026} (Fig. 1a). Each contains a main report body that describes the dairy size and operations, including herd characteristics (animal ages),  manure and wastewater generation, and farmland application. The reports also have attachments including manure and wastewater transport manifests, documenting off-site exports. Of the 795 annual reports analyzed, 375 contain a total of 1,196 transport manifests with 1,264 transport records. A single manifest can include multiple hauls, with unique dates, destinations, and tonnage. In total, the extraction yielded 828 manure transportation records and 435 wastewater transportation records (see SI Methods 0.1.1).

Annual reports and manifest attachments are extracted separately, each with an automated extraction followed by a manual review against the source PDF (Fig. S2). Annual reports for 580 dairies were parsed by Robertson et al 2026 \cite{Robertson_2026} from machine-readable copies of the standard Water Board form under Order R5-2007-0035 \cite{Dairy_Order_2007}. We parse a further 215 annual reports with an agentic approach that also reads scanned images and consultant-prepared reports with differing formats (Fig. 1b). Manifests were submitted by dairies using 14 different template variants (Fig. S1, Table S1), one of which is submitted via a software interface and the rest of which are filled out either by hand or with PDF editing tools. We use keyword-based text extraction based on form field titles to produce a summary file (Table S2) and one row per manifest. We then manually review and correct errors in the extracted information for each manifest (Fig. S3). Full extraction, OCR and review procedures are given in the SI Methods 0.1.1.

When expressing dairy facility size in animal units (AUs), as opposed to head count, we convert the herd classes based on values reported by Water Board in Order R5-2008-0066 \cite{Dairy_Order_2008}: milk cows at 1.40 AU, dry cows at 1.12 AU, large heifers at 1.02 AU, small heifers at 0.49 AU, and calves at 0.29 AU.

We further process several categories of manifest fields: haul size and quantity, haul dates, destination type, and origin and destination locations to determine the breakdown of endpoint types and timing variations for dairy manure and wastewater transportation. Hauls are binned as 10-ton or 20-ton loads \cite{Harrigan_1997} based on self-reported size or assumed from aggregate haul characteristics. Haul dates are taken from the first and last haul date provided, and tonnage is spread evenly across this date range (SI Methods 0.1.1). To process manifest destination, we use the manifest destinations categories included in the software-based manifest template: Farmer, Composting Facility (including manually-filled destinations naming a fertilizer company like ``Hyponex''), Broker, and Other (e.g. ``Garden'', ``Spreader''). We use a ``Blank'' category for missing destination type. We also characterize varying degrees of detail within the destination address type: street address, PO Box, place description, or no address at all.
 
For the manure manifests, we assume all transport events occur via truck. For wastewater transportation manifests, we distinguish transport mode as truck, pipeline, or unknown. There is no explicit field for the mode of wastewater transport so we determined the wastewater transportation mode based on keywords in the hauling company and ``method used'' fields (SI Table 3). A much larger proportion of the 435 wastewater records lack a usable destination type or transport mode, so we do not route these manifests, but incorporate them into analyses of data quality, destination types, and nitrogen accounting.

\subsection{Manure Nitrogen Accounting}
Dairy annual reports are designed to account for nutrients \cite{CVRWQCB_2013_Order}, rather than manure in itself, so we perform accounting on a nitrogen-basis. The total nitrogen available to a dairy comes from manure nitrogen after ammonia volatilization, which is assumed to be a flat 30\% in the annual report form. From this ``available nitrogen" we subtract the nitrogen that each dairy reportedly applies to its own land and exports. The residual nitrogen is labeled ``unaccounted-for" (Fig. 3a). For the Water Board-wide balance (Fig. 3a), we use a subset of 702 dairies with valid values for each of the 7 nitrogen accounting categories. Nitrogen concentrations in exported manure and wastewater are estimated using a range of conversion factors to capture uncertainty \cite{Davis_Tech_Guide} (see SI Methods 0.1.4). 

Most annual reports also state an overall nitrogen applied-to-removed ratio across each dairy's fields, defined as the ratio of nitrogen applied to the nitrogen removed by crops at harvest. The Dairy Order sets out a limit of 1.4 on a per-field basis. However, we use a whole-facility summary ratio as a proxy, as each dairy can have many fields, each with multiple crop rotations. We also extract each dairy's facility-wide crop nitrogen removal. We then create a metric "nitrogen pressure" for each dairy, a ratio defined as the available N divided by the implied crop-removed nitrogen. A nitrogen pressure above one means that the dairy generates more nitrogen than its land is taking up. We then explore whether nitrogen pressure and the overall applied-to-removed ratios correlate with a higher fraction of exports or unaccounted-for nitrogen (see SI Methods 0.1.2).

\subsection{Truck Routing and Emissions}

Each of the 828 manure records represents a unique origin-destination-pair and typically multiple trips, 150,333 10-ton hauls and 41,444 20-ton hauls in total. During pre-processing, we drop 78 pairs with missing destination information and failed geocodes, as well as those outside the water board study area or with assessor parcel number addresses that were geocoded outside of the origin dairy county (see Fig. S6). We then route the remaining pairs using a minimum travel time algorithm over an OpenStreetMaps road network \cite{OSM_California_2025}, clipped to a 25km buffer of Water Board boundary \cite{Water_Board_Boundaries}. We then manually review every origin, destination, and modeled path by hand, discarding routes with inaccurate geocodes or circuitous routing (See SI Methods 0.1.3, Fig. S3, Fig. S7). Trucking emissions and fuel usage rates are derived from CARB's EMFAC Model (v2.1.1) \cite{EMFAC_Emissions} (see SI Methods 0.1.3).

In total, 586 records are successfully routed and validated (70.8\%), accounting for 79,007 10-ton hauls, 29,439 20-ton hauls, and 1,378,209 tons of manure (59.1\%). We test whether excluded routes differ significantly from validated routes by their manure manifest tonnage and the origin dairy size (Fig. S7, Table S5).

To determine whether exposure to manure trucking differs across income groups, we compare the modeled road network against block group-scale median household income sourced from the 2020-2024 American Community Survey \cite{NHGIS_2026}. We buffer each road segment carrying manure trucking hauls by 500m to account for near-road pollutant exposure, approximating previous literature estimates of pollutant decay and health impacts based on distance from roadways \cite{Karner_2010, Zhu_2002, Gauderman_2007}. Truck-miles are spread uniformly across that buffer, allotted to block groups by the share of each buffer falling within them, then aggregated by median household income quintile over 4,851 block groups with available income statistics. We repeat this process to calculate exposure for Census-Designated Places (CDPs), towns, and cities from the 2024 Census Place dataset \cite{Census_Places, Census_Place_Shp_2024}.

\subsection{Manure Management Fraction}
CARB's GHG Emissions Inventory seeks to quantify the amount of methane emissions from dairy manure management by estimating the final manure management systems based on the fraction of total cattle head in each system. These systems include anaerobic digesters, anaerobic lagoons, land application, liquid/slurry, or solid storage. The Inventory then estimates the methane emissions inherent in each of these management streams \cite{CARB_2000_2023_GHG_Methods, CARB_2000-2023_GHG}. To compare what dairy annual reports and manifests record against the GHG Emissions Inventory assumptions, we estimate the fraction of generated manure ending in these manure management pipelines on the nitrogen basis defined in Section 2.2. See SI Methods 0.1.4 for more details.

\begin{figure}[H]
 \centering
        \includegraphics[width=1.0\textwidth]{Figures/Figure_2_Manifest_Characteristics.png}
 \caption{Large fractions of exported manure and wastewater have uncertain destination types and transportation methods. a) Exported manure tonnage (U.S. tons) by dairy-reported destination type reveals the plurality of destination types are left blank. b) Exported manure tonnage by destination address form. Most tonnage (71.9\%) went to a complete street address; 28.1\% went to a place description, a PO box, or to an unknown address. c) Approximately half of exported wastewater is sent to farmers, while the other half has no listed destination type. d) Across 435 wastewater records, the majority of exports have an uncertain transportation mode.}
\label{ManifestCharacteristics}
\end{figure}

\begin{figure}[H]
 \centering
        \includegraphics[width=1.0\textwidth]{Figures/Figure_3_Nitrogen_Accounting.png}
 \caption{Dairies with lower relative crop uptake capacity have larger nitrogen accounting issues. a) A waterfall chart showing the overall nitrogen balance across 702 Central Valley dairies which have complete data for each bar. Nitrogen generated from manure cannot be fully accounted for by ammonia volatilization, land applications, and exports. This accounting includes 129 dairies reporting more nitrogen applied and exported than generated. b) Facility-scale fractions of exported and unaccounted-for nitrogen as a share of post-ammonia loss dairy nitrogen (available nitrogen) are plotted against   pressure, the ratio of available nitrogen to nitrogen uptake from crops on dairy-owned fields. Lines represent a moving median with a window of 90 facilities, while the bands represent the interquartile range of dairies in each window. The typical dairy has a median post-ammonia loss unaccounted-for nitrogen fraction of 39.1\%. c, d) Spearmann rank correlations for dairy applied-to-removed ratio and nitrogen pressure against four outcomes. Bars are 95\% bootstrap confidence intervals over 2,000 draws. Nitrogen pressure shows a stronger relationship to unaccounted-for nitrogen, exported nitrogen, and herd size than the applied-to-removed ratio.}
\label{NitrogenAccounting}
\end{figure}

\section{Results}
\subsection{Characterizing Dairy Operational Data}

Reports cover a majority of dairy facilities and waste in California. Of 716 complete dairy annual reports, the total herd size was 1,783,114 head, which generated 30.7 million tons of fresh manure and 23.4 billion gallons of wastewater in 2024. From 828 manure transport records, we find that 2,331,698 tons of manure were trucked to 482 unique destinations. Wastewater exports totaled 894 million gallons. The reporting formats of annual reports and manifests are inconsistent. In total, we identify 4 types of annual reports and 14 different types of manure and wastewater transportation manifests (Table S1, Figure S1). Of the 795 analyzed annual reports, 690 (87\%) used the standard Water Board form while the remaining 105 used a consultant format (Fig. S1a). Of the analyzed dairy annual reports, 716 appeared complete in including both a herd size $>$1 and had complete information on generated manure and wastewater nitrogen. 

The operational records also suggest data integrity issues. First, the plurality of exported manure tonnage (38.4\%) did not report a destination type. The remainder listed farmers (25.6\%), brokers (20.2\%), composting facilities (13.9\%), or other (2.0\%) as destinations (Fig. 2a). In total, 654,920 tons of exported manure tonnage carried either no address or an incomplete address (such as a locality name, PO Box, or road name without a number),  28.1\% of total manure exports (Fig. 2b). This includes the single largest manure destination in the dataset, a PO Box for a produce company in Bakersfield that received 163,632 tons of manure across 18 records, 7.0\% of all manure exports (Table S6). Generally, facilities trucked lower tonnages of manure in winter months, and the highest tonnages were trucked in May and October (Fig. S4). Second, the 435 wastewater transportation records encompassed 894 million gallons of exports and had higher uncertainty in both destination and transportation mode. While 50.4\% of wastewater exports went to farmers, 49.5\% had blank destination types (Fig. 2c). Third, we were unable to determine the transportation mode for 73.0\% of exports; the remaining exports were split between pipelines (23.1\%) and tanker trucks (3.9\%) (Fig. 2d, SI Table 3).

%We also note the absence of manure digesters as a manure or wastewater end-destination. The earliest version of the manure manifest template was dated from the 2007 dairy general order R5-2007-0035 for the Water Board, prior to a surge of close to 100 manure digesters being built \cite{ERA_2024}. The subsequent updates to the template have remained largely unchanged. Thus, within the annual report and manifest datasets, the manure entering anaerobic digesters is not detectable.

\subsection{Manure Nitrogen Accounting}

We estimate that approximately 190,000 tons of nitrogen were generated from dairy manure and wastewater in 2024 (Fig. 3a). Of this, 57,000 tons volatilized as ammonia, 52,000 tons were applied to dairy-owned fields, and 60,000 tons were exported, leaving 21,000 tons of nitrogen unaccounted-for. Larger dairies had somewhat higher fractions of unaccounted-for nitrogen (Fig. S9d, Spearman $\rho$ = 0.185, p$\leq0.001$), but in much larger quantities (Fig. S9c, Spearman $\rho$ = 0.648, p$\leq0.001$). 

Of dairies reporting both an applied-to-removed ratio and crop nitrogen removal above zero (N=684), the median nitrogen pressure was 2.3, indicating that the median dairy generated more than twice the nitrogen it could remove via crop uptake. As expected, nitrogen pressure was associated with higher fractions of exported nitrogen (Fig. 3d, Spearman $\rho$ = 0.392, p$\leq0.001$). But consistent with the hypothesis that farms may experience cost pressures with transportation, the fraction of unaccounted-for nitrogen also increased with nitrogen pressure (Fig. 3d, Spearman $\rho$ = 0.410, p$\leq0.001$). The applied-to-removed ratio correlated with unaccounted-for nitrogen (Fig. 3c, Spearman $\rho$ = 0.139, p$\leq0.001$), though less than the nitrogen pressure metric.

Few dairies (23) self-reported a nitrogen applied-to-removed ratio over the Water Board threshold of 1.4 (Fig. S5). Dairies in the 1.2-1.4 ratio range had slightly larger fractions (median 0.448) of unaccounted-for nitrogen than other dairies (median 0.356) by Mann-Whitney test (p=0.005). In the counterfactual where all unaccounted-for nitrogen were assumed to be applied to fields, 457 of the 673 dairies that report crop nitrogen removal would exceed the 1.4 threshold, with a median applied-to-removed ratio of 1.95 (Fig. S5).

\subsection{Truck Routing and Emissions}

Of the 828 manure trucking records in the 2024 annual report dataset, 586 (70.8\%) were routed successfully and validated, accounting for 59.1\% of manifest tonnage. In total, we estimated the dataset contained 150,333 10-ton truck hauls and 41,444 20-ton truck hauls, of which 79,007 10-ton hauls and 29,439 20-ton hauls were routed and validated. Accounting for loaded outgoing trips and unloaded return trips, the successfully routed and validated manure trucking routes cover 2,425,147 miles annually, 1,763,807 miles by 10-ton capacity trucks and the remaining 661,340 miles by 20-ton capacity trucks. Using the EMFAC 2024 per-mile emissions rates for the San Joaquin Valley, we estimated PM$_{\text{2.5}}$ emissions of 127.0 kg, NOx emissions of 3,109.5 kg, and diesel fuel consumption of 377,795 gallons, resulting in 3,850 metric tons of CO$_{\text{2}}$ emitted \cite{EMFAC_Emissions}.

Manure hauls were generally local. Across validated routes, 89.2\% of trips were less than 25 miles (Fig. 4d). The longest trip (89.4 miles) was between Troost Dairy and a farmer in Tracy, San Joaquin County. The average route trucked was 10.1 miles one-way. The ten longest trucking routes contained 8.7\% of all truck miles traveled, but only 1.7\% of routed manure exports by ton (Table S7).

The routed manifests represent a subset of true routes, as 29.2\% of routes and 40.9\% of manure export tonnage excluded from the analysis. Manifests with larger tonnage were modestly more likely to be dropped overall  (Mann-Whitney p$\leq0.001$, r=0.19), from failed geocodes (Mann-Whitney p$\leq0.001$, r=0.45), and bad destination geocodes identified in review (Mann-Whitney p=0.001, r=0.22). The same pattern holds for manifests originating at larger dairies (Fig. S7, Table S5). As dropped manifests tended to be larger in tonnage, our routing analysis may miss some highly-impacted communities along these routes.

The most notable substantive finding from this routing analysis is that emissions fall disproportionately on more vulnerable communities. Across 4,851 block groups within the Water Board, exposure to manure trucking generally fell as median household income increased, on both a per-resident basis (Spearman $\rho$ = -0.182, p$\leq0.001$) and a per-mi$^2$ basis (Spearman $\rho$ = -0.194, p$\leq0.001$). The lowest income quintile block groups experience 0.40 truck-miles per person, compared to 0.12 truck-miles in the highest income quintile (Fig. 4e). The communities facing the highest total truck-miles were the cities of Tulare (50,664 mi), Visalia (33,804 mi), and Turlock (30,192 mi) (Table S8). The top 20 communities by per-resident exposure were all unincorporated CDPs, led by Dos Palos Y (32.17 truck-miles per resident). CDPs also make up 15 of the top 20 communities exposed per-mi$^2$, led by Lakeside CDP with 10,852 truck-miles per~mi$^2$ (Tables S9, S10). Exposure is concentrated, with the 5\% of residents with the highest per-capita truck-miles experiencing 46.4\% of the total exposed-place truck-miles.

\begin{figure}[H]
 \centering
        \includegraphics[width=1.0\textwidth]{Figures/Figure_4_Network_And_Exposure.png}
 \caption{Modeled manure trucking networks exposures are concentrated with low-income and unincorporated communities. a) A map of dairies in the Central Valley Water Board shown with County borders. Dairy icon size is scaled by the total animal unit size of the facilities in 2024 (see Section 2.1, SI Methods 0.1.1). These icons represent a subset of 686 dairies with complete herd table of dairy cows, dry cows, heifers, and calves. Dairy locations were geocoded from addresses provided in annual reports. b) Total 10-ton and 20-ton truck passes per road segment from modeled manure trucking routes. c) Census block group-scale median household income based on the 2020-2024 American Community Survey \cite{NHGIS_2026}. d) Distribution of one-way manure trucking route distances in 2024. e) Truck traffic exposure by median household income quintile, normalized by truck-mile exposure per-resident and per~mi$^2$. Error bars represent 95\% bootstrap confidence intervals. f) Truck-miles per resident across cities and Census-Designated Places (CDPs). Lower population CDPs tend to face higher per-resident exposure than towns or cities in the Central Valley (Spearman $\rho$ = -0.440, N=141, p$\leq0.001$).}
\label{RoutingAnalysis}
\end{figure}

\subsection{Manure Management Fraction}

The manure nitrogen endpoints we estimate from annual reports are on the same order as the CARB GHG Inventory's manure management endpoints for land application (17.4\% vs. 10.6\%), liquid/slurry (25.6\% vs. 20.2\%), and dry storage (5.5\% vs. 9.9\%) (Fig. 5b). However, the anaerobic digester and anaerobic lagoon manure fractions cannot be estimated from annual reports and manifests. Thus, the fractions of manure nitrogen going to these management processes is mixed within the unknown component, which encompasses manure and wastewater exports with unknown destinations and other unknown losses, such as potential groundwater infiltration. Additionally, we find that exported nitrogen reported in annual reports exceeds exports specifically documented in manifests by 36,079 tons. We estimate that 22-66\% of exports reported in these annual reports are accounted-for in manifests, depending on assumptions for exported manure and wastewater nitrogen content (see SI Methods 0.1.4). Determining final management methods for this un-manifested portion of nitrogen is not possible.

\section{Discussion}

Operational data for industrial dairies in California's Central Valley reveal missingness of data, but also trends in the unaccountability of nitrogen. The median dairy had 39.1\% of their post-ammonia loss nitrogen unaccounted-for by their annual reports (Fig. 3b), driven by larger dairies (Fig. S9), and large fractions of exported nitrogen reported in annual reports were not accounted for by transportation manifests (Fig. 5b, Section 3.2). Among the manifested exports, large fractions of exported manure and wastewater had no listed destination type (Fig. 2a,c) or usable destination addresses (Figs. 2b, S6). This limits the applicability of this regulatory data for a core part of the operative Dairy Order \cite{CVRWQCB_2013_Order}: accounting for where manure goes to ensure there is no land over-application.  

The Water Board sets a maximum applied-to-removed ratio for nitrogen and dairy facilities are penalized for exceedances. However, both the land application and crop removal values at a facility are self reported. Dairies with low capacity to remove their generated nitrogen via crop uptake (i.e. higher nitrogen pressure) face economic pressure, as transporting and exporting manure can be a large expense for dairies \cite{USDA_2003}. Many dairies in the Central Valley face this pressure; the median dairy here generated more than twice the nitrogen that it was able to land-apply in 2024. Dairies with higher nitrogen pressure had higher exports, but these did not close the accounting gap; unaccounted-for nitrogen also correlated strongly with nitrogen pressure (Section 3.2, Fig. 3b). One way to incentivize complete nitrogen accounting at dairies could be to automatically record any unaccounted-for nitrogen as land applied nitrogen.  This is what the Acid Rain Program does for imputing missing acid rain data \cite{Giles_2022}. Indeed, if we apply the counterfactual that all unaccounted-for manure was applied to fields, the majority of dairies would not meet the regulatory applied-to-removed ratio of 1.4 (Fig. S5, Section 3.2). 

In September 2026, the California State Water Resources Control Board (SWRCB) approved a new dairy order WQ 2026-0028, focusing on ongoing issues with dairy nitrogen reporting in the Central Valley Water Board \cite{SWRCB_2026_Order}. SWRCB cite a 2019 report \cite{Dairy_Monitoring_2019} which found that dairies impact regional groundwater quality, yet have ``a substantial amount of ``unaccounted-for'' manure nitrogen'', which they attribute partly to ``imprecision in measurement of applied nitrogen and irrigation water.'' The SWRCB proposed dairy order also highlights data quality issues and suggests requiring Water Board to report dairy land application data in a ``usable format'', as opposed to the existing and varied PDF reports \cite{SWRCB_2026_Order}. Although the 2013 Water Board General Dairy Order states that the ``general section of the annual report shall be completed on an annual report form provided by the Executive Officer'' \cite{CVRWQCB_2013_Order}, the 103 annual reports completed by Innovative Ag Services and Livingston Dairy Consulting do not utilize the required format. Additionally, the 2013 Dairy Order requires the use of tracking manifest form ``Attachment D of Order No. R5-2013-0122'' \cite{CVRWQCB_2013_Order}, but only a fraction of manifests use this form (SI Table 1). While agentic approaches can enable processing of non-standardized operational data, the Water Board could simplify facility-scale nutrient accounting and enforcement of nutrient land application requirements by enforcing existing reporting mandates. 
\begin{figure}[H]
 \centering
        \includegraphics[width=1.0\textwidth]{Figures/Figure_5_CARB_Manure_Management.png}
 \caption{Annual reports have comparable land application and liquid management fractions to the CARB greenhouse gas inventory, but a large fraction of manure nitrogen has an uncertain management endpoint. a) Dairy manure management system fractions used in CARB’s greenhouse gas inventory to estimate methane emissions based on manure treatment pathways. b) Fate of manure and wastewater nitrogen as estimated based on 2024 dairy annual reports and manure and wastewater transport manifests across 737 dairies that report both a herd table and nitrogen generation. The denominator for these nitrogen endpoint percentages is the nitrogen generated in dry manure after the 30\% ammonia volatilization loss, a total of 275 million lbs of nitrogen. Error bounds represent uncertainty in nitrogen contained in manure and wastewater exports \cite{Davis_Tech_Guide}. If a dairy reports exports, but no manifests are available or manifests do not account for all exports, this difference is assigned to the ``Unknown" column, as are exports with a blank destination type. The 16.5\% ``unaccounted-for” component is the difference between total nitrogen generated and all land application and export outflows. Anaerobic lagoon and digester fractions are unknown, and thus denoted by ``?''.}
\label{CARBCompare}
\end{figure}

CARB and its regional air boards, as parallel agencies to the California Water Boards, still rely on decades-old EPA survey information for dairy manure management fractions \cite{EPA_Inventory_Annexes_2013} for their methane calculations in the Greenhouse Gas Emissions Inventory \cite{CARB_2000-2023_GHG, CARB_2000_2023_GHG_Methods}. However, the California dairy landscape has been rapidly changing due to consolidation in the past two decades, with average herd size increasing by 43\% between just 2017 and 2022 \cite{ERA_2024}. We observe that the percent of manure nitrogen going to the land application and liquid management fractions are on the same order as CARB estimates, but that the manure fractions going to lagoons and digesters are not currently calculable from annual reports (Fig. ~\ref{CARBCompare}). Understanding whether dairies employ anaerobic lagoons or digesters is critical for methane estimation, as lagoons are highly methane-intensive \cite{CARB_2000-2023_GHG, Owen_2015, Naranjo_2020}, and leaking digesters may have significant methane emissions factors as well \cite{Wainer_2025}.

Our manure truck routing analysis estimated that manure trucks burned at least 377,795 gallons of diesel, releasing $\text{CO}_2$ equivalent to annual enteric emissions from 1,065 dairy cows \cite{CARB_2000-2023_GHG}. These modeled routes disproportionally flowed through lower-income block groups and rural, unincorporated Census-Designated Places (CDPs) (Fig. 4e,f). Lower income communities have previously been shown to face higher relative exposure to traffic \cite{Rowangould_2013, Gunier_2003}, which contributes to adverse respiratory health outcomes \cite{Gauderman_2007}. CDPs, which do not have municipal government representation or authority to regulate heavy vehicle traffic \cite{CA_Code_35700}, carry the highest per-resident and per-area exposures (Tables S9-10). Heavy trucking miles for dairy manure transportation are minor in comparison to total traffic flows, being approximately equal to the total daily miles driven within the City of Merced \cite{CDOT_2024}. Additionally, there is uncertainty regarding how well modeled trucking routes match true routes. Still, these trucking routes are a previously unquantified pollution source of intensive animal agriculture that affect low-income communities in the Central Valley.

Overall, we find that annual reports and manure manifests provide critical data for scoping environmental and equity implications of intensive dairy operations in California. Their utility is limited by addressable issues: lack of standardization, lack of information on key manure management methods, unaccounted-for nitrogen flows, and dairy noncompliance with export reporting requirements. SWRCB's recent dairy order underscores these issues, directing the Central Valley Water Board to implement a ``whole-farm'' nitrogen accounting strategy that ``requires all nitrogen at the dairy to be accounted for'' \cite{SWRCB_2026_Order}. They suggest required data may include amounts of manure treated by ``any manure management process'', which would theoretically encompass anaerobic lagoons and digesters. Such improvements to reporting would drastically increase the utility of annual reports for both pollution accounting and allow regulators to balance the benefits and externalities of intensive animal agriculture in California.

\ack{All network analysis and data visualization tasks were performed on Stanford's Sherlock Computing Cluster (RRID:SCR\_026161)}

%\funding{Sample text inserted for demonstration.}
% This section is a list of funder names and grant numbers

\roles{I. Avery Bick:
Conceptualization, Methodology, Formal analysis, Investigation, Visualization, Writing – original draft, Writing – review \& editing.
\newline
\newline
Daly Wettermark:
Formal analysis, Investigation, Writing – original draft, Writing – review \& editing.
\newline
\newline
Gina Hervey:
Conceptualization, Methodology, Writing – review \& editing.
\newline
\newline
Elena Eneva:
Supervision, Writing – review \& editing.
\newline
\newline
Daniel E. Ho:
Supervision, Writing – review \& editing.
}

% List author names and the contributions made to the article, using terms from the NISO Contributor Roles Taxonomy (CRediT) https://credit.niso.org

\section*{Data availability statement}
All annual reports, manifests, geospatial data, tabular data, and code required to reproduce results and figures in this paper are available at Zenodo.com/xxx.yyy.zzz.

% For more information on IOP Publishing's research data policy see: https://publishingsupport.iopscience.iop.org/questions/research-data/


\bibliographystyle{iopart-num}
\bibliography{refs}

\clearpage
\setcounter{section}{0}
\setcounter{subsection}{0}
\setcounter{subsubsection}{0}
\setcounter{figure}{0}
\renewcommand{\thefigure}{S\arabic{figure}}

\section*{Supplementary Information}
\subsection{Supplementary Methods}
\subsubsection{Characterizing Dairy Operational Data}
We extract information from annual reports and transportation manifests using a combination of methods. Of the 859 annual reports received by Robertson et al 2026 \cite{Robertson_2026}, 795 make up the analysis sample for this study. The other 64 were dropped due to having no recoverable text, being a non-dairy facility, or having reporting methods not comparable to the rest of the sample set, such as the Provost and Pritchard consulting firm which uses a different nitrogen reporting scheme. Robertson et al 2026 performed an earlier parse of 580 machine-readable and standard Water Board formatted annual reports under Order R5-2007-0035 \cite{Dairy_Order_2007}. Here, we parse an additional 215 annual reports using an agentic approach that allows for inclusion of scanned images, as well as consultant-derived reports of varying formats. All agent-extracted rows are reviewed by hand against the annual reports using an interactive tool available in this paper's associated Zenodo repository (Fig. S2). For plotting CAFO locations for visualization purposes (Fig. 4a), dairy addresses provided on facility annual reports are geocoded using the US Census Bureau Geocoding Service \cite{census_geocoder} and Nominatim \cite{nominatim} services. Of 795 dairies, 767 were geocoded.

To identify manure and wastewater manifest attachments within annual report PDFs, we first use a keyword search. We try two approaches to extract readable text from the PDF. We first use PyMuPDF to extract embedded text, when available. When embedded text results come up with fewer than 50 characters, we run optical character recognition (OCR) using Tesseract. For those pages, we render the images at 300 DPI using PyMuPDF. We run Tesseract with Page Segmentation Mode (PSM) 0 to identify the text orientation on the images, then the default PSM 3 to segment text regions of the document. From the extracted text, we use a list of terms that appear in the manifests, but nowhere else in the boilerplate text of the reports, to identify manifest-specific pages for extraction. The terms are ``hauler info'', ``destination'', ``method used'', ``operator shall'', ``d-2'',  ``solids content'', ``hauler signature'', ``hauling event'', and ``complete one''. We then use the search terms ``instructions'' and ``information'' to identify page 1 of a new manifest, as the first page of all templates included an instructions list. If the immediate following page also contains ``instructions'', we treat the identified manifest as a 1-page manifest. If not, we include the subsequent page assuming a 2-page manifest. For manifest pages that are deemed to have low-quality OCR outputs from Tesseract, based on having more than 15\% non-alphanumeric characters, we run the manifest pages through the LLMWhisperer API, using the free tier API key.

We use the python script to save each automatically-identified manifest to a standalone PDF as well as a .txt file. The .txt file includes page number delimiters to identify the original PDF page source. Following the automatic manifest identification, we manually screened the report PDFs to confirm the number of manifests present and add any missing manifests to the results manually. Our initial automatic count found 1,169 manifests. Manual screening identified in 2 false-positives (0.2\%) and 59 missed manifests (5.0\%) resulting from poor-quality, upside-down, or out-of-order scanned pages. After manually corrections, we retained 1,196 signed documents containing 1,264 manifest records: 828 manure and 435 wastewater records.

After identifying the manifests, we organize the data from each manifest’s extracted text. We again apply automatic extraction followed by manual verification. We found 14 distinct page formats used for the manifest attachments. The formats were between 1 and 3 pages and referenced various Central Valley Regional Water Board general order numbers, or no specific order. Three formats included a header identifying the consultant used to prepare the report and manifests. Though the exact language and format varied, the data fields included remained largely consistent. A summary of the templates identified is given in Table S1 and Figure S1b.

Of the 59 missed manifests, 56 (from 20 dairies) were missed by the keyword search due to poor scan quality. For these reports, we specified the manifest pages manually prior to OCR and extraction. One partial manifest only included the second page and was counted but does not contain a numerical record.

When grouping pages into manifests, we attach page 2 if they contains expected text such as ``amount hauled'', ``written agreement'', ``certification'', or a ``D-2'' page header. We stop attaching further pages once the manifest contains a ``certification'' section. We include a third page for forms labeled ``Page 2 of 3''.

We assign each manifest to a template (Table S1) using keyword rules applied to its extracted text. The rules are ordered so that consultant- and permit-specific forms (e.g. ``CDQAP'', ``CAG018001'') are matched before the general order forms. We identified the general order templates by comparing manifests against the blank Attachment D forms published with Orders R5-2007-0035 \cite{Dairy_Order_2007} and R5-2013-0122 \cite{CVRWQCB_2013_Order}, as well as Attachments F and G of Order R5-2010-0118 and Attachment D of Order R5-2017-0058. We distinguish templates with identical first pages using the order number on page 2 (e.g. R5-2007-0035 vs. R5-2013-0122) or differences in the written-agreement date (2007 vs. 2009). The one-page R5-2007-0035 Variant 5 template lists hauls in a table. We treat each table row as a separate record.

We remove 27 manifests that duplicate another manifest. 14 are exact copies. 13 are CAG018001 NPDES manifests (Form 4) that repeat an attachment for the same hauling event within the same report. In this case, we retain the Attachment D version which provides more detail. These 13 are counted in the template percentages in Table S1.

We extracted the following information from each manifest template: Origin Dairy Name, Origin Dairy Address, Hauler Name, Hauler Address, Destination Name, Destination Type, Destination Contact Address, Destination Address, Destination Nearest Cross Street, Destination Parcel Number, Haul Date, Amount Hauled, and Solids Percentage (for manure manifests), and Method Used to Determine Amount Hauled. 

We calculate accuracy of the automatically extracted value of each field against the manually validated value for 1,092 automatically-extracted records. Two values are treated as matching if the normalized string (ignoring case, punctuation, and spacing) of the two agree, with further normalization done to addresses using the libpostal library. Of the compared records, 620 (56.8\%) required at least one manual correction. 486 (44.5\%) required a correction to a field relevant to the main manuscript analyses: haul dates, addresses, parcel numbers, cross streets, hauled amounts, or haul count. Accuracy exceeded 90\% for ``Is Pipeline'', ``Total Manure Amount'', ``Last Haul Date''. Accuracy was lowest for the destination fields (``Destination Address'' 74.4\%, ``Destination Contact Address'' 58.9\%, ``Destination Assessor Parcel Number'' 40.5\%) due to formatting differences and handwriting (Fig. S10).

We clean the text for extra whitespace and common OCR errors (e.g. ``1BS'' rather than ``LBS''). To test whether a geocoded result is street-level, we parse the address returned by the geocoder with the libpostal library and require it to carry a house number or road cross-section. We discard results that geocode only to a city, county, or ZIP code level.

We categorize manure manifest destination addresses (Fig. 2b) based on the level of detail given. First, we check for empty addresses (``blank'' category). Then we check for "PO Box" strings. The ``Place Description'' category is assigned when a street name is present without a street number, or for general descriptions such as a ZIP code or town name. Lastly, the ``Street Address'' category is assigned for street names preceeded by a number, and for assessor parcel number addresses. 

We estimate trucking haul size in two bins: 10-ton and 20-ton hauls. Two facilities explicitly mention using these truck types. Other facilities reporting tonnage per haul tend to cluster around these values. Only a small fraction of manifests report the mass trucked per haul. Using the manifests that report haul size, we determine an approximate distribution of 10-ton and 20-ton hauls. We assign manifests reporting haul sizes of 5-15 tons and 15-25 tons to the 10- and 20-ton bins, respectively. The 10-ton bin represents 64\% of total hauled manure by mass among these manifests. For all remaining manifests, we apply this distribution to each manifest's hauled amount to estimate the portion carried by 10- or 20-ton trucks, then divide by trick capacity, rounding to the nearest whole haul count. Wastewater volume per manifest was calculated using several methods based on the information provided, including include flow over time and haul volume (Table S4).

To process haul dates, we use the last and, when available, the first haul date reported in the manifests. If a single haul date is reported, we attribute all hauling to that date. If a range of dates is provided, we split the total manure hauled evenly across the range of days provided. If no haul dates are provided, we do not include the data in haul timing analysis (Fig. S4). Across 828 manifest records, 817 contain a haul date and 567 carry only a last haul date.


\subsubsection{Manure Nitrogen Accounting}

When estimating per-facility unaccounted-for nitrogen (Fig. 3b-d), we start with manure nitrogen post-ammonia volatilization. Ammonia loss is estimated at 30\% of total generated manure nitrogen, the reduction that is applied within the operative dairy annual report \cite{CVRWQCB_2013_Order}. We define "available nitrogen" as manure-generated nitrogen post-ammonia loss. The fractions of unaccounted-for nitrogen shown in Fig. 3b-d are fractions of this available nitrogen that are unaccounted-for after accounting for nitrogen outputs to land application (manure and wastewater) and exports.

Robertson et al 2026 \cite{Robertson_2026} included wastewater as a distinct nitrogen source in their accounting of dairy nitrogen. A comment \cite{Dairy_Order_Comments_2026} submitted regarding the new SWRCB dairy order \cite{SWRCB_2026_Order} contended that nitrogen from wastewater was a subset of manure-generated nitrogen and not a distinct source. The annual reports are ambiguous regarding this distinction. Dairy water discharges can include nitrogen from manure, feed, milk, and other sources \cite{SWRCB_2026_Order, USDA_2008, Ali_2006}. As a conservative approach, we do not include reported wastewater nitrogen generation as a distinct nitrogen source for dairies, instead assuming that this nitrogen is a subset of the reported nitrogen generated from manure (see Fig. 3a). If wastewater nitrogen were included as a separate nitrogen source, the unaccounted-for nitrogen across the Water Board would increase by approximately 37,000 tons in 2024.

The applied-to-removed ratio for nitrogen is reported by dairies on a per-field basis and a facility-wide basis, which incorporates nitrogen applied and removed across all fields operated by the dairy. A per-field applied-to-removed nitrogen ratio of 1.4 is a compliance metric under the Water Board's R5-2013-0122 Order \cite{CVRWQCB_2013_Order}. As annual reports can contain many applied-to-removed ratios across fields and crop cycles, we use facility-wide applied-to-removed ratios as a proxy (Fig. 3c, Fig. S5). We extract facility-wide crop-removed nitrogen values from annual reports from the standard form's "Actual crop nutrient removal" field and the Innovative Ag Services form's "Crop Nitrogen Removal" field. The Livingston consulting firms do not report an aggregate crop nitrogen removal figure and are not accounted-for here. Nitrogen pressure is then calculated by dividing the available nitrogen by this crop-removed nitrogen. Nitrogen pressure is thus not defined for dairies when no nitrogen was land-applied or removed. To test how applied-to-removed ratios would change in the counterfactual scenario where unaccounted-for nitrogen was land-applied, we add each dairy's unaccounted-for nitrogen with its total reported land-applied nitrogen, then divide by the reported crop-removed nitrogen (Fig. S5). Dairies with no reported crop-removed nitrogen and incomplete nitrogen accounting information were excluded.

\subsubsection{Truck Routing and Emissions}
For facilities where the origin address is missing from the manure manifest, we backfill using the dairy address listed in the main report body. Next, we geocode the destination address where the manure or wastewater was delivered. As most manifest templates include multiple options to report destination address, we implement a priority order for geocoding: Destination Latitude/Longitude, Destination Parcel Number, Destination Nearest Cross Street, Destination Address, Destination Contact Address, and Hauler Address. Destination address parcel numbers are looked up via the California Department of Water Resources Statewide Parcel Geocoding Service\cite{dwr_parcel_geocoding}, keeping candidate parcels whose county matches the destination county reported on the manifest. The geocoding service uses different formats from county-to-county (e.g. ``025-200-003'' in Madera or ``012025013'' in Stanislaus). We first search each parcel number as written. If no parcel is returned for the search county, we retry its 9-digit form with and without dashes. We convert 16-digit parcel numbers with zeros to a 9-digit form (e.g. ``00025-0140-0005-0000'' to ``0025-140-005''). Destination cross streets are only used where two distinct street names are given. The “Destination Contact Address” is used if it appears to be a valid street address. Because this field is listed in the same sub-section as “Destination Address” on the manifest, we assume that the destination location and contact information are the same when the former is left blank. We use the “Hauler Address” if it contains the term “Farm”, “Compost,” or “Fertilizer.” We do not consider destinations with PO Boxes and addresses that only include a city or county without street information for the routing analysis, but these are maintained in the export accounting.

We use the Python library OSMNX (v2.1.0) \cite{OSMNX_2025} to preprocess the road network dataset: reprojecting the road network into UTM Zone 10N (EPSG:32610) to allow for distance calculations in meters, filling in missing edge lengths using planar geometry, and filling in missing road segment speed limits. Missing speed limit data are filled based on the road segment category (e.g. highways with a speed limit of 100kph, and residential streets with a speed limit of 30 kph). Any road segments missing both speed limit and road segment category information are assumed to have a speed limit of 40 kph. We then apply the NetworkX library (v3.6.1) \cite{NetworkX_2008} to optimize trucking routes with a goal of minimizing travel time between each origin-destination pair. To account for round trips, we assume that manure truck return paths retrace the outbound path from the origin dairy. Thus, each road segment carries every modeled haul twice. However, the return trip does not affect estimates of manure trucked per road segment. The total, validated one-way loaded travel distance is 1,212,574 miles, and including return trips, 2,425,147 miles.

Truck emissions and fuel usage rates are derived from CARB's EMFAC Model (v2.1.1) for the San Joaquin Valley Air Pollution Control District in 2024 \cite{EMFAC_Emissions}, which overlaps with large portions of the Water Board. We assign 10-ton truck trips to ``T7 Single Dump Class'' and 20-ton truck trips to  ``T7 Tractor Class 8''. The NOx exhaust rate for these classes is 1.21 g/mile and 1.47 g/mile for the 10-ton and 20-ton classes, respectively; $\text{CO}_2$ is emitted at 1,621.27 and 1,497.15 g/mile, respectively. $\text{PM}_{2.5}$ emissions factors sum truck exhaust with tire wear and brake wear, giving a total of 0.0516 g/mile for the 10-ton class and 0.0545 g/mile for the 20-ton class. Diesel consumption is derived from $\text{CO}_2$ emissions rates and a conversion factor of 10.19 kg $\text{CO}_2$ per gallon of diesel burned \cite{EIA_CO2_per_gal}, giving 6.29 and 6.81 miles per gallon. We calculate emissions from routed and validated truck miles traveled only, and do not incorporate other emission sources such as truck engine starting or idling, which also produce significant emissions.

\subsubsection{Manure Management Fractions}

To create our estimates of dairy manure management fractions on a nitrogen basis, we start with previously-validated 795 dairy annual reports, then drop reports for dairies with a herd population $\le$1, as well as 1 non-dairy operation. In total, we use a universe of 737 dairies, generating 275 million lbs of nitrogen after ammonia loss. We then link manure and wastewater transport manifests to these dairies to account for exports, with 778 of 828 manure records and 429 of 435 wastewater records accounted-for. On transport manifests, exports are listed as tonnage (for manure) and gallons (for wastewater). For manure manifests, we use self-reported solids percentage when available, or assume tonnage-weighted mean of 48.5\% when this value is not reported. To convert manure tonnage back to a nitrogen-basis, we use a distribution of values: minimum, median, and maximum of solid nitrogen content (1.2, 2.1, and 3.5\%) across 29 solid manure samples from the University of California Davis Cooperative Extension's Technical Guide \cite{Davis_Tech_Guide}. To convert wastewater exports to a nitrogen basis, we use Guide's minimum, median, and maximum of per-dairy median lagoon-water nitrogen concentrations (164, 353, 645 \textmu{}g/g)  \cite{Davis_Tech_Guide} , with a 1kg/L density of wastewater assumed. Any exports reported in annual reports, but not manifests, are listed in the ``unknown'' nitrogen destination column. However, this estimation has large error bounds due to uncertainty in exported manure and wastewater nitrogen concentrations.

Dairy annual reports detail manure and wastewater generation, nutrient applications, and exports, but do not detail usage of anaerobic lagoons and digesters. Thus, nitrogen flowing into these management fractions are encompassed in our `unknown'' nitrogen fraction. Additionally, it is possible that process wastewater in annual reports, which is applied to farmland or exported, may be digester digestate.

\begin{figure}
    \centering
    \includegraphics[width=1\linewidth]{Figures/Figure_S1_Reporting_Formats.png}
    \caption{Breakdown of a) dairy annual report and b) extracted manure and wastewater transport manifests in 2024 by format that are analyzed in this study. Manifest formats include those associated with dairy orders (R5-2007-0035 \cite{Dairy_Order_2007}, R5-2013-0122 \cite{CVRWQCB_2013_Order}, etc.) and custom formats.}
    \label{ReportingFormats}
\end{figure}

\begin{figure}
    \centering
    \includegraphics[width=1\linewidth]{Figures/Figure_S2_Annual_Report_Review_Example.png}
    \caption{Screenshot of the interactive HTML review application used to facilitate manual review of agent-derived Water Board dairy annual report data. Incorrect extracted values were updated manually on the right column. This process was completed for 215 annual reports that were not included in Robertson et al 2026 \cite{Robertson_2026} due to irregular formats.}
    \label{AnnualReportReviewTool}
\end{figure}

\begin{figure}
    \centering
    \includegraphics[width=1\linewidth]{Figures/Figure_S3_Geocoding_Routing_Review_Example.png}
    \caption{Interactive HTML review application of manure transport network analysis. The user reviews the origin dairy (left) and destination (middle) address and geocode, as well as the modeled route (right). Any origin-destination pairs with incorrect geocodes or routing were removed from further analysis.}
    \label{RoutingReviewTool}
\end{figure}

\begin{figure}
    \centering
    \includegraphics[width=1\linewidth]{Figures/Figure_S4_Manure_Trucked_Per_Month.png}
    \caption{Exported manure tonnage by month of haul in 2024. Manure trucking is higher during the warm growing season, and lower during the wet season, during which limitations on applications to saturated soils exist. \cite{CVRWQCB_2013_Order}}
    \label{}
\end{figure}

\begin{figure}
    \centering
    \includegraphics[width=1\linewidth]{Figures/Figure_S5_Applied_Ratio.png}
    \caption{Reported whole-farm nitrogen applied-to-removed ratios at Central Valley dairies in 2024 (black lines) versus counterfactual scenario where all unaccounted-for nitrogen is assumed to be land-applied (grey bars). The dotted line represents the 2013 Dairy Order threshold of 1.4 \cite{CVRWQCB_2013_Order}.  Only 23 dairies self-reported ratios higher than 1.4 in 2024. In the counterfactual scenario we see that 457 of the 673 dairies which report crop nitrogen removal would exceed the limit if unaccounted-for nitrogen were assumed to be land-applied.  In the reported applied-to-removed ratios, dairies in the threshold set (ratio of 1.2-1.4) are more likely to have a higher fraction of unaccounted-for nitrogen by Mann-Whitney test (p=0.001, rank-biserial=0.146).}
    \label{}
\end{figure}

\begin{figure}[H]
    \centering
    \includegraphics[width=1.0\textwidth]{Figures/Figure_S6_Manure_Sankey.png}
    \caption{Sankey diagram illustrating the data cleaning pipeline for manure transport manifests to prepare the trucking route networking analysis. Manifests were dropped when they: provided no destination, the destination locations were outside of the Water Board study region, the geocoding failed due to insufficient information, an Assessor Parcel Number (APN) address was geocoded outside of the dairy’s origin county, and when manual review determined that the origin or destination geocode or routing were incorrect. Of the 828 manure manifest records, 586 (70.8\%) were retained for the routing analysis.}
    \label{SankeyDiagram}
\end{figure}

\begin{figure}
    \centering
    \includegraphics[width=1\linewidth]{Figures/Figure_S7_Routing_Failure_Modes.png}
    \caption{Failure modes of manure manifests dropped from routing analysis (Fig.~\ref{SankeyDiagram}) by a-c) manure amount per manifest, and d-f) the herd size of the origin dairy. Overall, manifests with larger manure tonnage and those from larger origin dairies were more likely to be dropped in general and specifically for failed geocodes or bad destination geocodes.}
    \label{RoutingFailureModes}
\end{figure}

\begin{figure}
    \centering
    \includegraphics[width=1\linewidth]{Figures/Figure_S8_Routing_Zooms.png}
    \caption{Zoomed-in maps of modeled 2024 manure truck trips near a) Visalia, Pixley, and Delano, and b) Stockton, Merced, and Modesto.}
    \label{RoutingZooms}
\end{figure}

\begin{figure}
    \centering
    \includegraphics[width=1\linewidth]{Figures/Figure_S10_Manifest_Extraction_Accuracy.png}
    \caption{Accuracy of automatic manifest field extraction compared to manually validated values for 1,092 automatically-extracted manure and wastewater records. Bars show the number of non-blank values per field from manual validation (green) and automatic extraction (orange). Points show the percentage of values where the automatic extractio nagrees with the validated value.}
    \label{ExtractionAccuracy}
\end{figure}

\begin{figure}
    \centering
    \includegraphics[width=1\linewidth]{Figures/Figure_S9_Unaccounted_Nitrogen.png}
    \caption{Histograms of a) unaccounted-for nitrogen per-dairy and b) unaccounted-for nitrogen fraction per dairy, defined as the fraction of each dairy's post-ammonia loss manure-generated nitrogen accounted for by reported land applications and exports (``available'' nitrogen). Histograms correspond to the 672 dairy sample in Figure 3b. The median dairy had 83,582 lb of unaccounted-for nitrogen, about 39.1\% of available nitrogen. c) Unaccounted-for nitrogen and d) unaccounted-for nitrogen fraction versus dairy size in animal units across 605 dairies reporting herd size breakdowns. Lines are a moving median with a window of 90 facilities.}
    \label{RoutingZooms}
\end{figure}


Table S1: Breakdown of 2024 Water Board manure and wastewater transport manifests by form template used

\begin{longtable}{| >{\raggedright\arraybackslash}p{2.3cm} | >{\raggedright\arraybackslash}p{1.8cm} | >{\raggedright\arraybackslash}p{2cm} | >{\raggedright\arraybackslash}p{1.5cm} | >{\raggedright\arraybackslash}p{1.6cm} | >{\raggedright\arraybackslash}p{2.8cm} |}
\hline
Template ID & Attachment Letter & General Order Referenced & Page Count & Percent of All Submitted Manifests*& File example \\
\hline
R5-2007-0035 Software & D & R5-2007-0035 & 2 & 33.7\% & 13701 S Hwy 59 El Nido 2024 Dairy AR \\
\hline
R5-2007-0035 General Order & D & R5-2007-0035 & 2 & 42.0\% & 20000 Cedar Ave Tracy 2024 Dairy AR (manifests 1–2) \\
\hline
R5-2007-0035 Variant 1 & D & R5-2007-0035 & 2 & 0.5\% & 13775 Murphy Rd Escalon 2024 Dairy AR \\
\hline
R5-2007-0035 Variant 2 & D & R5-2007-0035 & 2 & 5.5\% & 348 E Monte Vista Rd Ceres 2024 Dairy AR \\
\hline
R5-2007-0035 Variant 3 & N/A & R5-2007-0035 & 2 & 0.1\% & 17399 E Avena Rd Escalon 2024 Dairy AR \\
\hline
R5-2007-0035 Variant 4 & D & R5-2007-0035 & 1 & 2.6\% & 2023AR\_F \& L Barcellos \#2\_7585 Avenue 152\_W Tulare \\
\hline
R5-2007-0035 Variant 5 & D & R5-2007-0035 & 1 & 2.2\% & 2024AR\_Faial Farms 2 Dairy \\
\hline
R5-2010-XXXX Attachment G & G & R5-2010-XXXX & 2 & 0.2\% & 2024AR\_The Dairy Inc\_6240 21st Ave\_Kings \\
\hline
R5-2010-XXXX Attachment F & F & R5-2010-XXXX & 2 & 0.1\% & 2024AR\_K A Vanderham Dairy\_12565 Ave 152\_Tulare \\
\hline
R5-2013-0122 General Order & D & R5-2013-0122 & 2 & 10.6\% & 13436 N Thornton Rd Lodi 2024 Dairy AR \\
\hline
R5-2017-0058 Variant 1 & D & R5-2017-0058 & 3 & 0.1\% & 2024AR\_Contente \& Co Ranch\_5730 20th Ave\_Kings \\
\hline
R5-2017-0000 &  & R5-2017-0000 & 2 & 0.6\% & 2024AR\_Cream Top Dairy\_13075 Ave 200\_Tulare \\
\hline
Provost Pritchard & N/A & N/A & 2 & 0.4\% & 2024AR\_Diamond D LLC Dairy\_9423 Idaho Ave\_Kings \\
\hline
UCANR & N/A & N/A & 1 & 0.1\% & 20000 Cedar Ave Tracy 2024 Dairy AR (manifest 3) \\
\hline
Unknown Variant 1 & N/A & N/A & 1 & 0.2\% & 2024AR\_Moonstone Dairy\_1305 Iona Avenue\_Kings \\
\hline
CAG018001 & 4 & CAG018001 & 1 & 1.2\% & 2024AR\_John Jongsma Dairy\_15434 Ave 192\_Tulare \\
\hline
Innovative Ag & N/A &  & 3 & 0.1\% & 2024AR\_South Creek Dairy\_11450 Avenue 64\_W Tulare \\
\hline

\end{longtable}
* The totals include the 1,196 manifests analyzed in the main manuscript plus 13 duplicates where a facility submitted two different form templates for the same hauling event(s). 

Table S2: All manifest parameter fields and post-processing applied

\begin{longtable}{| >{\raggedright\arraybackslash}p{4.5cm} | l | >{\raggedright\arraybackslash}p{2cm} | >{\raggedright\arraybackslash}p{4.5cm} |}
\hline
\textbf{Parameter Name} & \textbf{Data Type} & \textbf{Manifest Type} & \textbf{Post-Processing Applied} \\
\hline
Origin Dairy Name & text & both &  \\
\hline
Origin Dairy Address & text & both & Geocoding \\
\hline
Hauler Name & text & both &  \\
\hline
Hauler Address & text & both & Geocoding \\
\hline
Destination Name & text & both &  \\
\hline
Destination Contact Address & text & both & Geocoding \\
\hline
Destination Address & text & both & Geocoding, querying the address together with its cross street and county first, if given \\
\hline
Destination Nearest Cross Street & text & both & Geocoding, only where two distinct street names are given and the geocoder returns a street intersection \\
\hline
Destination Type & text & both & Standardization of categories \\
\hline
Destination Assessor Parcel Number & text & both & Geocoding against the statewide parcel locator, within the expected county \\
\hline
Destination Latitude & numeric & both &  \\
\hline
Destination Longitude & numeric & both &  \\
\hline
Haul Date & text & both & Parsed into First Haul Date (if applicable) and Last Haul Date \\
\hline
Is Pipeline & boolean & both & Set true if ``pipeline'' or ``ditch'' is included in the manifest text \\
\hline
Is Trucked & boolean & both & Always true for manure. Set true if number of hauls, trucks, or tanks is explicitly mentioned for wastewater. Otherwise ``False/Unsure'' \\
\hline
Total Process Wastewater Exports (Gallons) & numeric & wastewater &  \\
\hline
Total Manure Amount (tons) & numeric & manure &  \\
\hline
Total Manure Amount (cubic yards) & numeric & manure & Converted to tons by dividing by Manure Density (cubic yards per ton) \\
\hline
Manure Density & numeric & manure &  \\
\hline
Manure Solids (\%) & numeric & manure &  \\
\hline
Manure Moisture (\%) & numeric & manure & Converted to Solids \% \\
\hline
Method Used to Determine Volume of Manure & text & manure & Broken out into Manure Mass per Haul (tons), Manure Volume per Haul (cubic yards), and Manure Number of Hauls as relevant \\
\hline
Method Used to Determine Volume of Wastewater & text & wastewater & Broken out into Wastewater Volume per Haul (Gallons), Wastewater Number of Hauls, Wastewater Hours Pumped (Hours), and Wastewater Pumping Rate (GPM), as relevant \\
\hline

\end{longtable}

Table S3: Determining trucked versus pipeline wastewater transport

\begin{longtable}{| >{\raggedright\arraybackslash}p{3.5cm} | >{\raggedright\arraybackslash}p{2.8cm} | >{\raggedright\arraybackslash}p{6.5cm} |}
\hline
\textbf{Transport Type} & \textbf{Detection Keywords} & \textbf{Text Examples} \\
\hline
Pipeline & ``pipeline'' or ``ditch'' &

(via concrete ditch) \newline
Pipeline \newline
LUMAR PIPELINE \\ \hline Trucked & ``load'', ``haul'', ``tank'', or ``hauler'' (without ``apply'' or ``applied'') &
247 Loads @ 9500 Gallons Per Load. \newline
Haulers Calc \newline
Hauler's Calculations \newline
100 Loads With Slurry Tankers-8400 Gallons/Load \newline
Number Of Loads Multiplied By Tank Volume \\ \hline \end{longtable}



Table S4: Methods used to calculate wastewater volume

\begin{longtable}{| >{\raggedright\arraybackslash}p{3.5cm} | >{\raggedright\arraybackslash}p{2.8cm} | >{\raggedright\arraybackslash}p{6.5cm} |}
\hline
\textbf{Calculation Method} & \textbf{Data Extracted} & \textbf{Text Examples} \\
\hline
Flowrate-based & Hours pumped, pumping rate (GPM) &

Gpm X Runtime \newline
Run-Times Multiplied By The Gpm Of The Lagoon Pump. \newline
6 Hours @ 1000 Gpm \newline
20 Hours At 1500 Gallons Per Minute \newline
Hours Pumped X 1200 Gpm \newline
8 Hours Of Liquid Manure @ 1000 Gallons Per Minute \\ \hline Pond level-based & Pond level drop &
Pond Drop Calculation Based On 5 Feet Drop In First Quarter \newline
Pond Drop Calculations \newline
PD \\ \hline Haul-based & Number of hauls, volume per haul &
247 Loads @ 9500 Gallons Per Load. \newline
370 Loads X 10000 Gallons \newline
Number Of Loads Multiplied By Tank Volume \newline
100 Loads With Slurry Tankers-8400 Gallons/Load \newline
Slurry Tankers - 250 loads at 8,400 gallons per load. \newline
Haulers Load Count-17 Loads At 9500 Gallons Per Load. \newline
150 Loads At 9500 Gals/Load \newline
60 Loads Of Surry @ 9500 Gallons Per Load. \newline
10000 Gallons Per Load X 282 Loads \newline
@ 10000 Gallons Per Load \\ \hline Other/Unspecified & Varies &
Tank Capacity Multiplied By The Number Of Tanks Applied \newline
Precipitation Records \newline
Estimated Annual Total To Independent Farming Field. \newline
PO \\ \hline \end{longtable}


SI Table 5: Routed manure manifests versus dropped manure manifests. Rank-biserial correlations and p-values are from two-sided Mann-Whitney tests of dropped groups against the routed group (Fig. S7).

{\footnotesize
\begin{longtable}{| >{\raggedright\arraybackslash}p{4.2cm} | >{\raggedleft\arraybackslash}p{0.8cm} | >{\raggedleft\arraybackslash}p{1.5cm} | >{\raggedleft\arraybackslash}p{2.0cm} | >{\raggedleft\arraybackslash}p{2.4cm} | >{\raggedleft\arraybackslash}p{1.2cm} |}
\hline
\textbf{Group} & \textbf{$n$} & \textbf{Median} & \textbf{Median difference vs routed} & \textbf{Rank-biserial correlation} & \textbf{$p$} \\
\hline
\multicolumn{6}{|l|}{\textit{Manure per manifest record (tons)}} \\
\hline
Routed & 586 & 1,320.0 &  &  &  \\
\hline
No destination address & 12 & 1,825.9 & 505.9 & 0.20 & 0.228 \\
\hline
Outside water board & 4 & 2,575.0 & 1,255.0 & 0.17 & 0.548 \\
\hline
Failed geocode & 51 & 4,593.0 & 3,273.0 & 0.45 & $<$0.001 \\
\hline
Cross-county APN & 11 & 600.0 & $-$720.0 & $-$0.56 & 0.001 \\
\hline
Bad Origin Geocode & 17 & 2,000.0 & 680.0 & 0.34 & 0.017 \\
\hline
Bad Destination Geocode & 80 & 2,202.8 & 882.8 & 0.22 & 0.001 \\
\hline
Routing Incorrect & 67 & 1,440.0 & 120.0 & 0.04 & 0.563 \\
\hline
All dropped & 242 & 2,010.8 & 690.8 & 0.19 & $<$0.001 \\
\hline
\multicolumn{6}{|l|}{\textit{Herd size of the origin dairy (animal units)}} \\
\hline
Routed & 500 & 2,215.0 &  &  &  \\
\hline
No destination address & 7 & 2,712.5 & 497.5 & 0.26 & 0.243 \\
\hline
Outside water board & 3 & 4,990.6 & 2,775.6 & 0.39 & 0.248 \\
\hline
Failed geocode & 43 & 5,125.2 & 2,910.2 & 0.34 & $<$0.001 \\
\hline
Cross-county APN & 11 & 947.8 & $-$1,267.2 & $-$0.51 & 0.004 \\
\hline
Bad Origin Geocode & 17 & 2,384.7 & 169.7 & 0.09 & 0.551 \\
\hline
Bad Destination Geocode & 65 & 4,086.2 & 1,871.2 & 0.16 & 0.031 \\
\hline
Routing Incorrect & 56 & 2,481.3 & 266.3 & $-$0.03 & 0.761 \\
\hline
All dropped & 202 & 3,028.0 & 813.0 & 0.11 & 0.019 \\
\hline
\end{longtable}
}


SI Table 6: Largest receivers of dairy-exported manure across 828 manure manifest records. The tonnage percent are tons of total manifest manure exports from manifests with provided addresses.

{\setlength{\tabcolsep}{4pt}\footnotesize
\begin{longtable}{| >{\raggedleft\arraybackslash}p{0.8cm} | >{\raggedright\arraybackslash}p{3.6cm} | >{\raggedleft\arraybackslash}p{1.3cm} | >{\raggedleft\arraybackslash}p{1.4cm} | >{\raggedleft\arraybackslash}p{1.4cm} | >{\raggedleft\arraybackslash}p{1.4cm} | >{\raggedleft\arraybackslash}p{1.8cm} | >{\raggedright\arraybackslash}p{1.3cm} |}
\hline
\textbf{Rank} & \textbf{Destination}& \textbf{Tonnage}& \textbf{Manifests}& \textbf{Origin dairies} & \textbf{\% of tonnage} & \textbf{Cumulative \%} & \textbf{Address type}\\
\hline
1 & PO BOX 81498 BAKERSFIELD CA & 163,632 & 18 & 6 & 7.17 & 7.17 & PO box \\
\hline
2 & 2312 ALPINE MEADOWS CT WASCO CA & 137,324 & 10 & 3 & 6.02 & 13.19 & street \\
\hline
3 & 7559 AVE 152 TIPTON 93272 & 106,713 & 14 & 10 & 4.68 & 17.87 & street \\
\hline
4 & 13214 EBERLE RD BAKERSFIELD CA & 105,833 & 16 & 3 & 4.64 & 22.51 & street \\
\hline
5 & 14982 ROAD 152 TIPTON CA 93272 & 60,138 & 8 & 5 & 2.64 & 25.15 & street \\
\hline
6 & PO BOX 27 EDISON CA & 53,574 & 6 & 2 & 2.35 & 27.50 & PO box \\
\hline
7 & 6219 E PANAM LANE BAKERSFIELD CA & 49,017 & 7 & 1 & 2.15 & 29.65 & street \\
\hline
8 & 10909 GLEN CLOE PLACE LAMONT CA & 47,987 & 7 & 1 & 2.10 & 31.75 & street \\
\hline
9 & EL NIDO 95317 COYOTE RD MERCED & 37,044 & 5 & 2 & 1.62 & 33.37 & street \\
\hline
10 & 4322 SACO RD BAKERSFIELD CA & 36,100 & 6 & 2 & 1.58 & 34.96 & street \\
\hline
11 & PO BOX 20370 BAKERSFIELD CA & 30,984 & 5 & 1 & 1.36 & 36.31 & PO box \\
\hline
12 & 19085 W BRADBURY TURLOCK 95380 & 28,427 & 12 & 1 & 1.25 & 37.56 & street \\
\hline
13 & SOUTH LAKE ROAD BAKERSFIELD CA & 27,638 & 3 & 1 & 1.21 & 38.77 & street \\
\hline
14 & 2312 ALPINE MEADOWS CT WASCO CA 93280 & 24,632 & 2 & 2 & 1.08 & 39.85 & street \\
\hline
15 & WILLIAMS FAMILY DAIRY COMPOST 6801 AVENUE 120 PIXLEY CA 93256 & 24,510 & 3 & 1 & 1.07 & 40.93 & street \\
\hline
16 & 19075 W THIRD AVE STEVINSON CA 95374 & 23,500 & 2 & 2 & 1.03 & 41.96 & street \\
\hline
17 & 9442 BEAR MTN BLVD 93311 & 21,225 & 2 & 2 & 0.93 & 42.89 & street \\
\hline
18 & 13076 AVE 368 VISALIA 93292 & 21,200 & 3 & 3 & 0.93 & 43.82 & street \\
\hline
19 & 7792 NICKLE RD DOS PALOS 93620 MERCED & 20,096 & 9 & 3 & 0.88 & 44.70 & street \\
\hline
20 & 23234 SECOND AVE STEVINSON 95374 & 17,820 & 5 & 3 & 0.78 & 45.48 & street \\
\hline
\end{longtable}
}


SI Table 7: The ten longest validated manure trucking routes sorted by one-way routed travel distance.


{\setlength{\tabcolsep}{2pt}\scriptsize
\begin{longtable}{| >{\raggedleft\arraybackslash}p{0.7cm} | >{\raggedleft\arraybackslash}p{0.8cm} | >{\raggedright\arraybackslash}p{1.1cm} | >{\raggedright\arraybackslash}p{1.45cm} | >{\raggedright\arraybackslash}p{1.0cm} | >{\raggedright\arraybackslash}p{1.8cm} | >{\raggedright\arraybackslash}p{1.45cm} | >{\raggedright\arraybackslash}p{1.05cm} | >{\raggedleft\arraybackslash}p{1.2cm} | >{\raggedleft\arraybackslash}p{0.8cm} | >{\raggedleft\arraybackslash}p{1.0cm} | >{\raggedleft\arraybackslash}p{1.27cm} |}
\hline
\textbf{Rank} & \textbf{Route ID} & \textbf{Origin dairy} & \textbf{Origin address} & \textbf{Origin county} & \textbf{Destination address} & \textbf{Dest. type} & \textbf{Dest. county} & \textbf{One-way miles} & \textbf{Total trips} & \textbf{Manure (tons)} & \textbf{Loaded truck-mi} \\
\hline
1 & 131 & Troost Dairy & 27000 Road 9 Chowchilla Madera 93610 & Madera & 20177 S Macaurthur Dr Tracy 95304 San Joaquin & Farmer & San Joaquin & 89.4 & 102 & 1,304.3 & 9,116 \\
\hline
2 & 132 & Troost Dairy & 27000 Road 9 Chowchilla Madera 93610 & Madera & 20177 S Macaurthur Dr Tracy 95304 San Joaquin & Farmer & San Joaquin & 89.4 & 43 & 554.0 & 3,843 \\
\hline
3 & 133 & Troost Dairy & 27000 Road 9 Chowchilla Madera 93610 & Madera & 20177 S Macaurthur Dr Tracy 95304 San Joaquin & Farmer & San Joaquin & 89.4 & 69 & 884.6 & 6,166 \\
\hline
4 & 375 & Flint Dairy & 6511 Flint Ave. Hanford Kings 93230 & Kings & 7745 E Sandymush Merced 95341 Merced & Farmer & Merced & 77.9 & 309 & 3,970.8 & 24,060 \\
\hline
5 & 74 & Van Tol Dairy No. 2 & 6756 County Road 30 Orland Glenn 95963 & Glenn & 82 West Liberty Road Gridle 95948 County Butte & Composting Facility & Butte & 64.7 & 53 & 685.0 & 3,430 \\
\hline
6 & 425 & Morning Star Dairy Inc & 10262 W Elkhorn Ave Burrel Fresno 93607 & Fresno & 18919 Ave 23 Chowchilla 93610 Madera Madera & Composting Facility & Madera & 50.7 & 154 & 1,968.2 & 7,811 \\
\hline
7 & 346 & Banos Dairy & 4667 Avenue 23 1/2 Chowchilla Madera 93610 & Madera & 1275 North Avenue Gustine 95322 North Avenue Merced & Broker & Merced & 47.5 & 140 & 1,800.0 & 6,650 \\
\hline
8 & 347 & Banos Dairy & 4667 Avenue 23 1/2 Chowchilla Madera 93610 & Madera & 1275 North Avenue Gustine 95322 North Avenue Merced & Broker & Merced & 47.5 & 164 & 2,100.0 & 7,790 \\
\hline
9 & 669 & Avalon Dairy & 27225 Rond Rd Waslo 93280 & Kern & 9442 Bear Mtn Blvd 93311 & Blank & Kern & 47.2 & 719 & 9,240.0 & 33,920 \\
\hline
10 & 383 & South Lest Dairy (Former Milk Time Dairy) & 12519 Road 17 Madera Madera 93637 & Madera & 7792 Nickle Rd Dos Palos 93620 Merced & Farmer & Merced & 46.3 & 59 & 750.0 & 2,731 \\
\hline
\end{longtable}
}



SI Table 8: Communities with largest total manure trucking miles, based on borders provided by the Census Place dataset \cite{Census_Place_Shp_2024}. Truck-miles are allotted based on the fraction of 500 meter road buffers falling inside the borders of each Census Place.

{\footnotesize
\begin{longtable}{| >{\raggedleft\arraybackslash}p{1.2cm} | >{\raggedright\arraybackslash}p{6.0cm} | >{\raggedleft\arraybackslash}p{3.5cm} |}
\hline
\textbf{Rank} & \textbf{Place} & \textbf{Truck-miles} \\
\hline
1 & Tulare city & 50,664 \\
\hline
2 & Visalia city & 33,804 \\
\hline
3 & Turlock city & 30,192 \\
\hline
4 & Bakersfield city & 27,371 \\
\hline
5 & Shafter city & 26,560 \\
\hline
6 & Hilmar-Irwin CDP & 20,097 \\
\hline
7 & Lakeside CDP & 19,585 \\
\hline
8 & Hanford city & 16,564 \\
\hline
9 & Tipton CDP & 16,065 \\
\hline
10 & Merced city & 14,572 \\
\hline
11 & Manteca city & 14,451 \\
\hline
12 & Modesto city & 13,468 \\
\hline
13 & Fresno city & 11,831 \\
\hline
14 & Pixley CDP & 9,935 \\
\hline
15 & Madera city & 9,187 \\
\hline
16 & Dos Palos Y CDP & 8,781 \\
\hline
17 & Rosedale CDP & 7,840 \\
\hline
18 & McFarland city & 7,263 \\
\hline
19 & Riverdale CDP & 6,583 \\
\hline
20 & Atwater city & 6,461 \\
\hline
\end{longtable}
}


SI Table 9: Communities with the most manure trucking miles per resident. Population is derived from 2020-2024 ACS five-year estimates \cite{Census_Places}.

{\footnotesize
\begin{longtable}{| >{\raggedleft\arraybackslash}p{1.2cm} | >{\raggedright\arraybackslash}p{6.0cm} | >{\raggedleft\arraybackslash}p{3.5cm} |}
\hline
\textbf{Rank} & \textbf{Place} & \textbf{Truck-miles per resident} \\
\hline
1 & Dos Palos Y CDP & 32.17 \\
\hline
2 & Lakeside CDP & 24.27 \\
\hline
3 & Stevinson CDP & 22.26 \\
\hline
4 & Tuttle CDP & 20.96 \\
\hline
5 & El Nido CDP & 10.49 \\
\hline
6 & Raisin City CDP & 7.61 \\
\hline
7 & Walnut Grove CDP & 6.41 \\
\hline
8 & Lanare CDP & 5.76 \\
\hline
9 & Tipton CDP & 5.63 \\
\hline
10 & Fairmead CDP & 4.56 \\
\hline
11 & Volta CDP & 4.18 \\
\hline
12 & Bowles CDP & 3.73 \\
\hline
13 & East Tulare Villa CDP & 3.46 \\
\hline
14 & Hilmar-Irwin CDP & 3.28 \\
\hline
15 & Traver CDP & 2.94 \\
\hline
16 & Pixley CDP & 2.66 \\
\hline
17 & Monmouth CDP & 2.24 \\
\hline
18 & Cherokee Strip CDP & 2.21 \\
\hline
19 & Artois CDP & 2.17 \\
\hline
20 & West Park CDP & 2.13 \\
\hline
\end{longtable}
}


SI Table 10: Communities with the most manure trucking miles per square mile. Area of communities is derived from Census Place outlines \cite{Census_Place_Shp_2024}.

{\footnotesize
\begin{longtable}{| >{\raggedleft\arraybackslash}p{1.2cm} | >{\raggedright\arraybackslash}p{6.0cm} | >{\raggedleft\arraybackslash}p{3.5cm} |}
\hline
\textbf{Rank} & \textbf{Place} & \textbf{Truck-miles per mi$^2$} \\
\hline
1 & Lakeside CDP & 10,852 \\
\hline
2 & Tipton CDP & 9,306 \\
\hline
3 & East Tulare Villa CDP & 8,444 \\
\hline
4 & Dos Palos Y CDP & 5,610 \\
\hline
5 & Hilmar-Irwin CDP & 5,117 \\
\hline
6 & Empire CDP & 4,262 \\
\hline
7 & Cherokee Strip CDP & 2,876 \\
\hline
8 & Stevinson CDP & 2,836 \\
\hline
9 & Tulare city & 2,443 \\
\hline
10 & Bear Creek CDP & 2,332 \\
\hline
11 & Patterson Tract CDP & 2,223 \\
\hline
12 & Escalon city & 2,157 \\
\hline
13 & Gustine city & 1,918 \\
\hline
14 & Pixley CDP & 1,892 \\
\hline
15 & Franklin CDP & 1,849 \\
\hline
16 & Turlock city & 1,798 \\
\hline
17 & Hughson city & 1,766 \\
\hline
18 & Raisin City CDP & 1,703 \\
\hline
19 & Riverdale CDP & 1,677 \\
\hline
20 & Traver CDP & 1,675 \\
\hline
\end{longtable}
}



\end{document}


